%\resizebox{\textwidth}{!}{%
\begin{tikzpicture}

% Central node
\node[dia/rbox=4cm, minimum height=1.5cm] (center) at (0,0)
  {\textbf{Педагогічна \\ефективність ІОР}};

% Pentagon vertices (5 characteristics with theoretical anchors)
% Top
\node[dia/rbox=4cm, minimum height=1.8cm] (int) at (90:5cm)
  {\textbf{Інтерактивність}\\[2pt]\scriptsize Dalgarno \& Lee:\\взаємодія учня};

% Upper-left
\node[dia/rbox=3.5cm, minimum height=1.8cm] (imm) at (162:5cm)
  {\textbf{Занурення / присутність}\\[2pt]\scriptsize Slater; CAMIL:\\психологічна присутність};

% Lower-left
\node[dia/rbox=4.5cm, minimum height=1.8cm] (ctx) at (234:5cm)
  {\textbf{Контекстуальність}\\[2pt]\scriptsize Ситуативне пізнання;\\експерієнціальне навчання};

% Lower-right
\node[dia/rbox=4.5cm, minimum height=1.8cm] (vis) at (306:5cm)
  {\textbf{Візуалізація}\\[2pt]\scriptsize CTML; Mulders et al.:\\репрезентаційна точність};

% Upper-right
\node[dia/rbox=3.5cm, minimum height=1.8cm] (adp) at (18:5cm)
  {\textbf{Адаптивність}\\[2pt]\scriptsize CLT; Poupard et al.:\\когнітивне навантаження};

% Arrows from vertices to center
\draw[dia/arrow] (int) -- (center);
\draw[dia/arrow] (imm) -- (center);
\draw[dia/arrow] (ctx) -- (center);
\draw[dia/arrow] (vis) -- (center);
\draw[dia/arrow] (adp) -- (center);

% Pentagon edges (dashed)
\draw[dia/line, dashed] (int) -- (imm);
\draw[dia/line, dashed] (imm) -- (ctx);
\draw[dia/line, dashed] (ctx) -- (vis);
\draw[dia/line, dashed] (vis) -- (adp);
\draw[dia/line, dashed] (adp) -- (int);

\end{tikzpicture}%
%}
